% Chapter 7 worked example. Included by ch07_resources.tex.
\section{Worked Example: Remote Monitoring: Economic Evaluation, Affordability, and HTA}
\label{sec:comp-ch07-worked-example}

\subsection{Purpose and Status}
This worked example demonstrates how the methods in Chapter~7 can be
translated into a reviewable institutional decision. All numerical inputs are
synthetic unless a source is identified explicitly. They are intended to show the
logic of the analysis, not to report empirical findings or establish a universal
threshold. Local use requires replacement of the assumptions with current,
versioned evidence and approval through the relevant clinical, managerial,
economic, legal, technical, and governance processes.

A public payer and hospital network are assessing an AI-supported remote-monitoring programme for people with heart failure. The programme may reduce emergency visits and admissions, but it requires devices, connectivity, monitoring staff, escalation capacity, and recurring platform fees. 

\subsection{Decision Problem and Comparators}
\textbf{Decision problem.} Is the programme cost-effective, affordable, and suitable for local adoption, and which evidence gaps justify a staged decision?

The comparator set is deliberately broader than the existing pathway. This prevents
a new technology from receiving a lower evidentiary threshold than staffing,
workflow, shared-service, conventional digital, or no-change alternatives.

\begin{longtable}{@{}p{0.08\textwidth}p{0.84\textwidth}@{}}
\caption{Comparators for the Chapter 7 worked example}
\label{tab:comp-ch07-comparators}\\
\toprule
\textbf{Option} & \textbf{Comparator or strategy} \\
\midrule
\endfirsthead
\toprule
\textbf{Option} & \textbf{Comparator or strategy} \\
\midrule
\endhead
\bottomrule
\endlastfoot
1 & Usual scheduled follow-up \\
2 & Non-AI remote monitoring with threshold rules \\
3 & AI-supported remote monitoring with nurse review \\
4 & Enhanced community nursing without remote monitoring \\
\end{longtable}

\subsection{Model and Decision Structure}
The analytical structure should follow the decision rather than the available
software. The team should map the causal pathway from intervention input to
information, human decision, action, outcome, resource consequence, and review.
The structure should also identify capacity constraints, uncertainty, subgroup
variation, implementation requirements, and events that would change the
recommendation.

A compact quantitative relationship used in this example is:
\begin{equation}
\mathrm{NMB}=\lambda\Delta E-\Delta C
\label{eq:comp-ch07-key-relationship}
\end{equation}
The equation is an organizing aid rather than a substitute for disaggregated evidence,
qualitative findings, or mandatory safeguards. Its variables and interpretation must
be defined in the local protocol.

\subsection{Synthetic Parameters and Evidence Sources}
Table~\ref{tab:comp-ch07-parameters} provides the base-case inputs. A
production analysis should add parameter distributions, correlations, structural
alternatives, data dates, system versions, and evidence-confidence ratings.

\begin{longtable}{@{}p{0.32\textwidth}p{0.22\textwidth}p{0.38\textwidth}@{}}
\caption{Synthetic parameters for the Chapter 7 worked example}
\label{tab:comp-ch07-parameters}\\
\toprule
\textbf{Parameter} & \textbf{Base value} & \textbf{Source or interpretation} \\
\midrule
\endfirsthead
\toprule
\textbf{Parameter} & \textbf{Base value} & \textbf{Source or interpretation} \\
\midrule
\endhead
\bottomrule
\endlastfoot
Eligible population & 8,000 & Commissioning estimate \\
Year-1 uptake & 35\% & Implementation scenario \\
Incremental cost per participant & EUR 1,180/year & Programme costing \\
Expected admissions avoided & 0.10/person-year & Evidence synthesis \\
QALY gain & 0.028/person-year & Economic model \\
Cost per admission & EUR 5,600 & Payer data \\
Implementation cost & EUR 1.4 million & Network estimate \\
Decision threshold & EUR 35,000/QALY & Illustrative local threshold \\
\end{longtable}

\subsection{Step-by-Step Analysis}
The sequence below is intended to make the analysis reproducible and to ensure
that evidence, economics, governance, and implementation develop together.

\begin{longtable}{@{}p{0.07\textwidth}p{0.55\textwidth}p{0.30\textwidth}@{}}
\caption{Analytical sequence}
\label{tab:comp-ch07-analysis-steps}\\
\toprule
\textbf{Step} & \textbf{Required work} & \textbf{Decision record} \\
\midrule
\endfirsthead
\toprule
\textbf{Step} & \textbf{Required work} & \textbf{Decision record} \\
\midrule
\endhead
\bottomrule
\endlastfoot
1 & Specify decision-maker, perspective, comparator, population, horizon, and intended scale. & Evidence and responsible owner to be recorded \\
2 & Estimate incremental costs, effects, QALYs, and net benefit. & Evidence and responsible owner to be recorded \\
3 & Present annual and cumulative budget impact under uptake scenarios. & Evidence and responsible owner to be recorded \\
4 & Separate economic cost from cash flow and productive capacity from cash-releasing savings. & Evidence and responsible owner to be recorded \\
5 & Assess subgroup access, connectivity, digital support, and distribution of opportunity cost. & Evidence and responsible owner to be recorded \\
6 & Identify whether remaining uncertainty should lead to adoption, a restricted pilot, price negotiation, or further research. & Evidence and responsible owner to be recorded \\
\end{longtable}

\subsection{Illustrative Outputs}
The outputs in Table~\ref{tab:comp-ch07-outputs} show how technical,
clinical, operational, economic, equity, and governance findings should be kept
visible rather than compressed into a single favourable or unfavourable score.

\begin{longtable}{@{}p{0.30\textwidth}p{0.24\textwidth}p{0.38\textwidth}@{}}
\caption{Illustrative model and decision outputs}
\label{tab:comp-ch07-outputs}\\
\toprule
\textbf{Output} & \textbf{Result} & \textbf{Interpretation} \\
\midrule
\endfirsthead
\toprule
\textbf{Output} & \textbf{Result} & \textbf{Interpretation} \\
\midrule
\endhead
\bottomrule
\endlastfoot
Incremental programme cost before offsets & EUR 9.44 million/year at full scale & Excludes initial implementation \\
Expected admission offset & EUR 4.48 million/year at full scale & Dependent on causal effectiveness \\
Illustrative ICER & Approximately EUR 22,900/QALY & Favourable at stated threshold \\
Year-1 budget impact & Positive and substantial & May remain unaffordable despite cost-effectiveness \\
Preferred decision & Phased adoption with negotiated price and access support & HTA and affordability separated \\
\end{longtable}

\subsection{Sensitivity, Uncertainty, and Subgroups}
At minimum, the completed analysis should vary the parameters most likely to
change the decision, test at least one credible alternative model structure, and
report subgroup results separately where performance, access, response, burden,
or opportunity cost may differ. Deterministic analysis should identify thresholds;
probabilistic analysis should be used where credible parameter distributions and
correlations are available. Scenario analysis is preferable when structural or
future uncertainty cannot be represented defensibly by one probability model.

The record should distinguish uncertainty that can be reduced through evidence,
uncertainty that can be managed through safeguards, and uncertainty that remains
too consequential to accept. Missing evidence should not be represented as an
average or neutral value without justification.

\subsection{Interpretation and Recommendation}
Support phased adoption where the comparative clinical evidence is credible, but link scale to annual affordability, observed service capacity, subgroup access, and updated budget impact. The payer and provider should align payment and implementation responsibilities.

The recommendation is conditional rather than permanent. It applies only to the
specified population, setting, intervention configuration, evidence cut-off, and
version. The following conditions should be incorporated into the approval,
contract, implementation, monitoring, or renewal record:

\begin{itemize}
 \item the exact intervention bundle and comparator are defined.
 \item annual and cumulative budget impact remain within an agreed envelope.
 \item non-digital access and connectivity support are funded.
 \item local causal and implementation evidence is collected.
 \item price, uptake, and service-capacity thresholds are included in the decision record.
\end{itemize}

\subsection{Decision Statement}
\begin{quote}
For the specified decision problem and comparator set, the institution should
\emph{[proceed / proceed with conditions / redesign / pilot / restrict / defer /
replace / retire]}. The decision applies to \emph{[population, setting, version,
and evidence period]}, is owned by \emph{[accountable authority]}, and will be
reviewed on \emph{[date]} or earlier if \emph{[trigger events]} occur. The
evidence, assumptions, unresolved uncertainty, mandatory safeguards, monitoring
thresholds, and exit arrangements are recorded in the accompanying dossier.
\end{quote}
